\section{引言}

城市群配送正同时经历网络协同与车队电动化.
多车场共享客户和车辆能够减少重复运输,
电动车也为降低道路直接排放提供了新的选择;
但二者并不会自然转化为可执行的低碳方案.
车辆在何处补能、何时充电、由哪个车场承担客户以及订单变化后如何继承既有任务,
会同时改变成本、排放和成员收益.
因此, 研究重点已不只是寻找更短的路线,
而是使协同路线在实体车辆、能源补给、成员参与和动态执行层面保持一致.

现有研究首先从不同侧面推进了这一问题.
在多车场协同方面,
Wang等\cite{ref:89}将资源共享纳入带时间窗的时变多车场路径优化,
Li等\cite{ref:93}进一步研究异构车队下的多车场电动车协同.
在分散决策条件下,
Shi等\cite{ref:90}同时考虑混合车队、成员参与和分时电价,
表明协同边界会受到运营主体和能源价格共同影响.
这些研究说明多车场资源共享不能只按系统总成本评价,
还应辨别客户责任、车场能力和参与条件对协同空间的约束.

混合车队和充电物理构成第二类关键问题.
传统燃油车与电动车在载重、续驶能力、补能时间和排放核算上存在本质差异,
简单按里程比例估算会改变路线排序甚至产生在完整模型下不可执行的方案.
Goeke和Schneider\cite{ref:27}采用与载重和行驶状态相关的能耗函数研究混合车队路径,
Wang等\cite{ref:80}讨论异构电动车队与非线性充电,
Nafstad等\cite{ref:94}则从精确算法角度表明异构充电技术和非线性补能会影响
时间与电量状态的联动可行性.
与此同时, 电力排放和电价均随时段变化,
仅比较车辆类型而忽略充电时刻, 难以判断电动化带来的真实减排.

实体车多趟与动态订单进一步增加了计划和执行之间的差距.
Zhen等\cite{ref:44}将配送趟生成与实体车衔接分别处理,
Zhao等\cite{ref:91}以动态规划分割和混合遗传搜索求解多趟时变路径,
说明“客户访问序列”和“实体车辆日内排班”需要分层处理.
动态路径研究则强调新增、取消和改量订单的实时响应\cite{ref:14,ref:18},
近期研究已把动态需求扩展到多车场电动车场景\cite{ref:92}.
然而, 若滚动优化只更新未服务客户而不冻结已执行路径、
车辆当前位置、电量、已发生充电和累计成本,
新方案就可能撤回历史决策或重复使用同一实体车.

上述进展分别覆盖了协同、混合车队、充电和动态调度,
但仍缺少一个贯穿“计划形成—完整核算—动态继承”的统一框架.
其核心难点不是简单叠加约束,
而是保证同一方案在四个层面同时成立:
路线层面覆盖客户并满足时间窗,
排班层面由有限实体车完成多趟衔接,
能源层面在具体日期、城市和半小时时段完成电价与碳强度映射,
组织层面不突破成员参与底线,
并在动态事件后继承已经发生的状态.
只有这些层面共同闭合, 系统成本和排放的改善才具有运营含义.

这一统一核算也对求解方法提出要求.
混合遗传搜索已在多车场、多属性和多趟时变路径中表现出较强的通用性
\cite{ref:76,ref:77,ref:91},
事件图与动态规划等问题结构也可嵌入遗传搜索以提高复杂约束处理能力
\cite{ref:95}.
但在本文问题中, 每次局部移动都立即完成实体车排班、非线性补能和分时排放核算,
计算代价过高; 若始终用同一简化代价搜索,
又可能遗漏在完整目标下更有价值的路线结构.
因此需要将快速代理搜索与统一完整核算分开,
同时保留不同代理视角产生的互补候选,
再用受时限约束的数学规划完成路线组合和最终择优.

基于上述问题, 本文研究时变碳强度下多车场混合车队动态协同配送,
主要工作如下.
首先, 建立配送趟—实体车排班两层模型,
把有限车队、多趟衔接、趟间充电、客户责任、成员参与和滚动状态继承
纳入同一可执行方案, 并统一核算运营成本与充电相关排放.
其次, 在固定路线和实体车排班下单独优化充电开始时刻,
从而把充电择时的边际效应与路线、车型变化区分开.
再次, 设计两阶段多视角混合遗传搜索—限时MIP路线池重组算法,
以燃油、简化电动和机制感知三种代理视角生成候选,
由完整模型统一检查, 通过同视角精英延续和两次限时路线组合提高候选覆盖,
最终保留完整模型成本最低的可行解.
最后, 在公开标准算例和中国三大城市群81个构造算例上检验算法,
并围绕协同、充电、参与和动态订单开展机制分析.

为明确本文相对既有研究的要素边界,
表\ref{tab:literature-comparison}从多车场、混合车队、实体车多趟、
动态事件、分时充电排放和成员参与六个方面进行比较.
